% answers_short.tex —— 1958 年·北京卷 · 简短答案速查
% 由 tools/make_short_answers.py 自 answers.tex 的 \gksolhead/\gkanswer 提取，
% 并与 1958/1958_北京_数学真题_答案要点.md 逐题比对后生成，勿手改（改源 MD 后重跑）。

\begin{gknotice}
\textbf{使用说明}：本节只给最终答案便于快速核对，完整推导见解析版。答案为编写者独立推导，\textbf{非官方标准答案}。
\end{gknotice}

\gkpart{题目}
\begin{itemize}[leftmargin=2.6em]
  \item[\textbf{1.}] $x=9$
  \item[\textbf{2.}] $20$ 位数
  \item[\textbf{3.}] $19600$ 种
  \item[\textbf{4.}] $\dfrac{22}{63}$
  \item[\textbf{5.}] $\sqrt{2}$
  \item[\textbf{6.}] 当两不同点非对径时能作且只能作 $1$ 个大圆；当两点是对径点时能作无穷多个
  \item[\textbf{7.}] $2a$
  \item[\textbf{8.}] (1) $y>0$ 当且仅当 $-1<x<3$；$y<0$ 当且仅当 $x<-1$ 或 $x>3$；(2) 顶点 $(1,4)$，对称轴 $x=1$；(3) 开口向下的抛物线，过 $(-1,0)$、$(0,3)$、$(1,4)$、$(3,0)$，画法见解析
  \item[\textbf{9.}] $\alpha$，$\beta$ 分别为 $25^\circ$ 与 $65^\circ$（可互换）
  \item[\textbf{10.}] $1$ 寸
  \item[\textbf{11.}] 见解析版
\end{itemize}
